\documentclass[11pt]{article}
\usepackage[margin=1in]{geometry}
\usepackage{amsmath,amssymb,amsthm,mathtools}
\usepackage[colorlinks=true,linkcolor=blue,urlcolor=blue]{hyperref}
\newtheorem{theorem}{Theorem}
\newtheorem{lemma}[theorem]{Lemma}
\newtheorem{corollary}[theorem]{Corollary}
\newtheorem{assumption}[theorem]{Assumption}
\newtheorem{remark}[theorem]{Remark}
\newcommand{\N}{\mathcal N}
\newcommand{\A}{\mathcal A}
\newcommand{\E}{\mathbb E}
\newcommand{\BMO}{\mathrm{BMO}}
\newcommand{\Carl}{\mathrm{Carl}}
\newcommand{\R}{\mathbb R}
\newcommand{\norm}[1]{\left\|#1\right\|}
\newcommand{\fint}{\mathop{\textstyle\int\kern-1.1ex-}\nolimits}
\title{A conditional noncommutative Poisson--Carleson decomposition\\
via stopped random walks}
\author{}
\date{October 21, 2026}
\begin{document}
\maketitle

\noindent
This note continues Remark \texttt{cd Mei}(ii) of
\texttt{nc-Davis\_v3.tex}. Its only nonstandard input is the following
endpoint Davis estimate. The argument below does not prove that input.
In particular, this is a consequence conditional on the stated version
of the manuscript's main theorem, rather than an independent proof of it.

\begin{assumption}[Davis estimate for regular filtrations]\label{ass:D}
There is a constant \(D_4\), independent of the finite probability
von Neumann algebra and the length of the filtration, such that
\[
 \norm{x}_{H_1}\le D_4
 \norm{(\mathcal E_nx)_n}_{L_1(\A;\ell_\infty)}
\]
for every \(L_2\)-martingale associated with a \(4\)-regular filtration.
The initial algebra is allowed to be nontrivial. We index its
conditional expectation by \(\mathcal E_0\); this is an actual
conditional expectation, not the zero map. The Hardy norm includes
the initial martingale term \(\mathcal E_0x\).
\end{assumption}

Let \((\N,\varphi)\) be a finite von Neumann algebra with normalized
faithful normal trace. Put
\[
 \norm{f}_{\BMO_{cr}}
 =\max\left\{
 \sup_I\norm{\fint_I(f-f_I)^*(f-f_I)\,dx}_{\N}^{1/2},
 \sup_I\norm{\fint_I(f-f_I)(f-f_I)^*\,dx}_{\N}^{1/2}
 \right\}.
\]
We use Mei's operator-valued BMO space, with bounded operator-valued
interval averages, modulo constants. For a positive operator-valued
Borel measure \(\nu\) on \(\R^2_+\), weak-star countably additive and
finite on bounded rectangles, set
\[
 \norm{\nu}_{\Carl}
 =\sup_I |I|^{-1}\norm{\nu(I\times(0,|I|])}_{\N}.
\]
An order bound \(\nu(T(I))\le A|I|1_\N\) is equivalent to this bound.
For a difference \(\mu=\mu_+-\mu_-\) of positive measures, its Poisson
balayage is
\[
 S_P\mu(x)=\int_{\R^2_+}P_y(x-t)\,d\mu(t,y),
 \qquad P_y(s)=\frac1\pi\frac{y}{s^2+y^2}.
\]

\begin{theorem}\label{thm:main}
Assume Assumption~\ref{ass:D}. If \(f=f^*\in\BMO_{cr}(\R,\N)\) has compact
support, there are a selfadjoint \(g\in L_\infty(\R)\bar\otimes\N\) and
finite positive \(\N\)-valued measures \(\mu_+,\mu_-\), supported in a bounded
rectangle of the upper half-plane, such that
\[
 \boxed{f=g+S_P(\mu_+-\mu_-)}
\]
and
\[
 \boxed{\norm{g}_\infty+\norm{\mu_++\mu_-}_{\Carl}
       \le C(D_4)\norm{f}_{\BMO_{cr}}.}
\]
The equality holds in \(L_1(\R\bar\otimes\N)\), and in particular locally
in \(L_1\) and distributionally. The constant is independent of the
matrix dimension of \(\N\). For a general complex \(f\), four positive
measures suffice.
\end{theorem}

\section*{1. A signed martingale representation}
For the filtration indexed by \(n\ge0\), use the inhomogeneous column
martingale BMO norm
\[
 \max\left\{
 \norm{\mathcal E_0(F^*F)}_\infty^{1/2},\,
 \sup_{n\ge1}\norm{\mathcal E_n
   |F-\mathcal E_{n-1}F|^2}_\infty^{1/2}\right\},
\]
and take the maximum with the row norm. This is the manuscript's
convention after shifting the first index from \(1\) to \(0\).
\begin{lemma}\label{lem:representation}
Under Assumption~\ref{ass:D}, if \(F=F^*\in\BMO(\A)\), for a finite
probability algebra \(\A\) with a \(4\)-regular filtration, then
\[
 F=G+\sum_{n\ge0}\mathcal E_n(v_n^+-v_n^-),
 \qquad G=G^*,\quad v_n^\pm\ge0,
\]
where
\[
 \norm{G}_\infty+
 \norm{\sum_n(v_n^++v_n^-)}_\infty
 \le C(D_4)\norm{F}_{\BMO}.
\]
Both positive series after conditional expectation converge in \(L_1\).
\end{lemma}

\begin{proof}
For \(x\in L_2(\A)\), define
\[
 Jx=(x,\mathcal E_0x-x,\mathcal E_1x-x,\ldots).
\]
The \(L_2\)-Doob inequality implies that \(Jx\in L_1(\A;c_0)\):
indeed, for \(x_m=\mathcal E_mx\), the maximal \(L_1\)-norm of the tail
\((\mathcal E_nx-x)_{n\ge m}\) is at most
\(C\norm{x-x_m}_2\), since the trace is finite and normalized.
Also
\[
 \norm{(\mathcal E_nx)_n}_{L_1(\ell_\infty)}
 \le2\norm{Jx}_{L_1(\ell_\infty)}.
\]
The \(H_1\)--BMO duality and Assumption~\ref{ass:D} therefore show that
\(Jx\mapsto\tau(Fx)\) is bounded by
\(C D_4\norm{F}_{\BMO}\).
Hahn--Banach and
\[
 L_1(\A;c_0)^*=L_\infty(\A;\ell_1)
\]
give a sequence \((w,u_0,u_1,\ldots)\) of that norm with
\[
 \tau(Fx)=\tau(w^*x)
       +\sum_n\tau\bigl(u_n^*(\mathcal E_nx-x)\bigr).
\]
The \(c_0\) duality is the endpoint assertion in Junge--Xu,
\emph{Noncommutative maximal ergodic theorems}, Remark 2.2(ii).
After replacing the representing sequence by its selfadjoint part, write
\[
 F=G+\sum_n\mathcal E_nu_n,\qquad
 G=w-\sum_nu_n.
\]
The latter sum converges weak-star with bounded norm.

For completeness, positivity with a bounded sum follows from the
factorization defining \(L_\infty(\ell_1)\), not from termwise spectral
positive parts. Choose balanced factors
\[
 u_n=\sum_k a_{kn}^*b_{kn},\qquad
 \norm{\sum_{k,n}a_{kn}^*a_{kn}}_\infty,\
 \norm{\sum_{k,n}b_{kn}^*b_{kn}}_\infty\le C D_4\norm{F}_{\BMO}.
\]
Polarization gives
\[
 u_n=\sum_{r=0}^3i^{-r}v_n^{(r)},\qquad
 v_n^{(r)}=\frac14\sum_k
 (a_{kn}+i^rb_{kn})^*(a_{kn}+i^rb_{kn})\ge0.
\]
Each \(\sum_n v_n^{(r)}\) is bounded by the same constant up to an
absolute factor. Taking real parts gives \(u_n=v_n^{(0)}-v_n^{(2)}\).
These positive sequences have bounded sums; their conditional
expectations have summable traces, proving \(L_1\) convergence.
\end{proof}

\section*{2. Scalar discrete harmonic kernels}
Take \(h=1/(2m)\). Let a simple symmetric nearest-neighbor random walk
on \(h\mathbb Z^2\) move in the four coordinate directions with probability
\(1/4\). Its horizontal starting point may be any real number.
Let \(Q^h_{jh}\) be the horizontal exit-distribution operator when the
walk starts at height \(jh\) and is stopped on height \(0\). Thus
\[
 Q^h_{jh}a(x)=\sum_{k\in\mathbb Z}p_j(k)a(x+kh),\qquad Q^h_0=I.
\]
The kernels are symmetric probability distributions and
\(Q^h_{jh}Q^h_{lh}=Q^h_{(j+l)h}\).

\begin{lemma}\label{lem:kernel}
There is an absolute \(C\) such that, for \(j\ge1\),
\[
 p_j(k)\le C\frac{j}{j^2+k^2}.
\]
For every Schwartz function \(\psi\),
\[
 \sup_{\substack{0\le y\le1\\y\in h\mathbb Z}}\,
 \sup_x |Q^h_y\psi(x)-P_y\psi(x)|\longrightarrow0.
\]
Moreover, if \(I\) has length \(L\), \(t\in I\), and \(0\le y\le L\),
then an unrestricted walk started at \((t,y)\) exits the lower boundary
inside \(DI\) with probability at least \(1/2\), for an absolute
dilation \(D\).
\end{lemma}

\begin{proof}
The horizontal Fourier multiplier is
\[
 \sum_kp_j(k)e^{ik\theta}
 =\exp[-j a(\theta)],\qquad
 a(\theta)=\operatorname{arcosh}(2-\cos\theta).
\]
This follows from the discrete harmonic equation
\(r+r^{-1}=4-2\cos\theta\), choosing the root \(0<r\le1\).
On \([-\pi,\pi]\), \(a(\theta)\simeq|\theta|\); its first derivative is
bounded on each side of zero and its second derivative is bounded there.
The integral formula gives \(p_j(k)\le C/j\). Integrating by parts
twice on the two sides of zero gives \(p_j(k)\le Cj/k^2\):
the jump of the first derivative contributes \(O(j)\), and the
\(L_1\)-norm of the second derivative of \(e^{-ja}\) is \(O(j)\).
Together these imply the displayed kernel bound.

Extend \(a\) periodically with period \(2\pi\).
Writing \(a_h(\xi)=h^{-1}a(h\xi)\), the multiplier of \(Q^h_y\) is
\(e^{-ya_h(\xi)}\), and \(a_h(\xi)\to|\xi|\) for fixed \(\xi\).
Fourier inversion and dominated convergence yield uniform convergence
for \(0\le y\le1\).
Finally
\(\sum_{|k|>Aj}p_j(k)\le C/A\); choose \(A\) large and then \(D\)
large enough. The assertion for \(y=0\) is immediate.
\end{proof}

\section*{3. Uniform BMO control of the stopped lift}
Translations and dilations reduce the theorem to
\(\operatorname{supp}f\subset[-1/8,1/8]\).
Write \(K=\norm{f}_{\BMO_{cr}}\); the case \(K=0\) is trivial.
Averages over \([-1/2,1/2]\), using that \(f\) vanishes on a fixed
fraction of this interval, give the order bounds
\[
 \int_\R f^*f\,dx\le C K^2 1_\N,\qquad
 \int_\R ff^*\,dx\le C K^2 1_\N.                 \tag{11}
\]
Indeed, if \(Q\) is this interval and
\(|\operatorname{supp}f|\le |Q|/2\), operator Cauchy--Schwarz gives
\(|f_Q|^2\le\frac12\fint_Q f^*f\); combine this with
\(\fint_Q|f-f_Q|^2\le K^2 1\). The row estimate is identical.

Set \(f_\epsilon=P_\epsilon f\).
It is bounded and operator-norm Lipschitz, its BMO norm is at most \(K\),
and it satisfies the same integrated square bounds (11).
We use the elementary Poisson-variance estimate
\[
 \sup_{x,y}\norm{
 P_y\bigl(|f_\epsilon-P_yf_\epsilon(x)|^2\bigr)(x)}_\N
 \le C K^2.                                      \tag{12}
\]
This follows by decomposing the Poisson integral into concentric dyadic
intervals: the square oscillation on the \(r\)-th interval relative
to the smallest interval average is at most \(C(r+1)^2K^2\),
and the Poisson weights decrease as \(2^{-r}\). The mean
\(P_yf_\epsilon(x)\) minimizes this operator-valued variance.
The same estimate holds for adjoints.

Choose \(h\) small enough that
\[
 h\,\operatorname{Lip}(f_\epsilon)\le K.
\]
Comparing \(p_j(k)\) with the Poisson mass of the interval of length
\(h\) centered at \(kh\), and using the Lipschitz bound, yields
\[
 Q^h_y|f_\epsilon-c|^2(x)
 \le C P_y|f_\epsilon-c|^2(x)
       +C h^2\operatorname{Lip}(f_\epsilon)^2 1_\N.       \tag{13}
\]
In particular (12)--(13) imply
\[
 \norm{Q^h_y|f_\epsilon|^2(x)
            -|Q^h_yf_\epsilon(x)|^2}_\N\le C K^2.       \tag{14}
\]

Start the walk at \((X_0,1/2)\), with \(X_0\) uniform on \([-R,R]\).
Let \(\tau\) be its first exit from \(0<y<1\), and thereafter keep
the position \(Z_n=(X_n,Y_n)\) fixed. The filtration reveals \(X_0\)
initially and then all four-way coins, including after \(\tau\).
It is \(4\)-regular. For a boundary function \(a\), define
\[
 T_h a=
 \begin{cases}
 a(X_\tau),&Y_\tau=0,\\
 Q^h_1a(X_\tau),&Y_\tau=1.
 \end{cases}
\]
Discrete harmonicity and optional stopping give exactly
\[
 \E_n(T_h a)=Q^h_{Y_n}a(X_n).                    \tag{15}
\]
Put \(F=T_h f_\epsilon\).
Its conditional future variance is bounded by (14): it is the stopped
conditional expectation of the unrestricted lower-exit terminal value
\(f_\epsilon(\xi)\).
Thus its little BMO square norm is \(O(K^2)\).
Regularity controls the big BMO terms by four times these terms.
The initial term is controlled by
\[
 \E_0|F|^2\le Q^h_{1/2}|f_\epsilon|^2(X_0)\le C K^2 1_\N,
\]
using (11), Jensen for \(P_\epsilon\), and (13).
Consequently
\[
 \norm{F}_{\BMO}\le C K                         \tag{16}
\]
uniformly in \(R,\epsilon,h\) subject to the preceding small-mesh condition.

\section*{4. Transfer of the representation and its Carleson bound}
Apply Lemma~\ref{lem:representation} to \(F\):
\[
 F=G+\sum_n\E_n(v_n^+-v_n^-),\qquad
 \norm{G}_\infty\le A,\quad
 0\le\sum_n(v_n^++v_n^-)\le A1_\N,\quad A=C(D_4)K.
\]
Define positive measures on the closed strip by
\[
 \mu_{R,h}^\pm(B)
 =2R\,\E\sum_n v_n^\pm\,\mathbf1_{\{Z_n\in B\}},
 \qquad B\subset\R\times[0,1].                  \tag{17}
\]
The variables \(v_n^\pm\) need not be adapted; only their bounded
positive sum is used.

For every interval \(I\),
\[
 (\mu_{R,h}^++\mu_{R,h}^-)(I\times[0,|I|])
 \le C A |I|1_\N.                              \tag{18}
\]
To see this, let \(\sigma\) be the first visit of the stopped path to
that box. If \(\sigma<\infty\), it is at or before strip exit.
From \(Z_\sigma\), continue an unrestricted walk to the lower boundary.
By Lemma~\ref{lem:kernel}, its exit lies in \(DI\) with conditional
probability at least \(1/2\). Hence
\[
 2R\,\mathbb P(\sigma<\infty)
 \le 2(2R)\mathbb P(\xi\in DI)\le2D|I|.
\]
The last inequality holds because convolution of the initial density
\(\mathbf1_{[-R,R]}\) with any probability kernel is at most \(1\).
Equation (18) now follows from the pointwise positive-sum bound.

Test the martingale representation against \(T_h\psi\), and use (15):
\[
 2R\,\E\bigl[(T_h f_\epsilon)^*T_h\psi\bigr]
 =\int_\R g_{R,h}(x)^*\psi(x)\,dx
   +\int Q^h_y\psi(t)\,d(\mu_{R,h}^+-\mu_{R,h}^-)(t,y).       \tag{19}
\]
This is an operator-valued weak identity, or equivalently can be tested
against every normal functional on \(\N\). The trace notation in
deriving it is suppressed in (19).
The function \(g_{R,h}\) is the adjoint image of \(G\).
It satisfies \(\norm{g_{R,h}}_\infty\le A\):
under Lebesgue-distributed initial horizontal positions the two strip
exit marginals are each \(\frac12 dx\); truncating initial positions
only decreases them. The top marginal is then convolved with the
probability kernel \(Q^h_1\).

Let \(R\to\infty\). The two exit marginals increase to
\(\frac12 dx\). Symmetry and the vertical semigroup property give
\[
 \lim_{R\to\infty}2R\,\E\bigl[(T_h f_\epsilon)^*T_h\psi\bigr]
 =\frac12\int_\R(f_\epsilon+Q^h_2f_\epsilon)^*\psi\,dx.       \tag{20}
\]
This extra \(Q^h_2f_\epsilon\) term is essential.

The bounds (18) give local weak-star compactness of the positive
operator-valued measures; bounded \(g_{R,h}\) have weak-star compactness
in \(L_\infty\). Subnets may be used if the predual is nonseparable.
Pass first in \(R\), and then take \(\epsilon\downarrow0\), \(h\downarrow0\)
with \(h\operatorname{Lip}(f_\epsilon)\le K\).
The kernel convergence in Lemma~\ref{lem:kernel}, together with
\(f_\epsilon\to f\) and \(Q^h_2f_\epsilon\to P_2f\) in \(L_2\), yields
\[
 \frac12(f+P_2f)=g_0+S_P(\lambda_+-\lambda_-).              \tag{21}
\]
At height zero the Poisson operator in this formula is understood
as the identity.

Here the limiting positive measures are supported in the closed strip
and still satisfy (18). For clarity, passage from local weak-star
convergence to (21) has no missing tail condition: for a compactly
supported smooth test function \(\psi\),
\[
 |Q^h_y\psi(t)|+|P_y\psi(t)|
 \le C_\psi(1+t^2)^{-1}\quad (|t|\text{ large},\ 0\le y\le1),
\]
uniformly in small \(h\). The measure of each horizontal annulus
of radius \(r\ge1\), including all heights, is \(O(Ar)\).
Thus its contribution to the pairing is \(O(A/r)\), summable over
dyadic annuli. On compact sets kernel convergence is uniform.
Open-box approximation passes (18) to the limit.

The height-zero parts obey
\(\lambda_\pm(B\times\{0\})\le C A|B|1_\N\).
They have weak-star measurable densities in
\(L_\infty(\R)\bar\otimes\N\) of norm \(O(A)\), and are absorbed into
\(g_0\). Also \(\norm{P_2f}_\infty\le C K\), by (11).
Multiplying (21) by two and absorbing these terms proves
\[
 f=g_1+S_P(\mu_+-\mu_-),\qquad
 \norm{g_1}_\infty+\norm{\mu_++\mu_-}_{\Carl}\le C(D_4)K,    \tag{22}
\]
with measures supported in \(0<y\le1\). At this stage the measures need
not have finite total mass.

\section*{5. Spatial localization and finite mass}
Choose \(0\le\chi\le1\), smooth and compactly supported, equal to \(1\)
on the support of \(f\), with absolute bounds on \(\chi,\chi'\)
in the normalized coordinates. If \(\mu=\mu_+-\mu_-\) is as in (22), put
\[
 E_\chi(x)
 =S_P(\chi\mu)(x)-\chi(x)S_P\mu(x)
 =\int P_y(x-t)(\chi(t)-\chi(x))\,d\mu(t,y).
\]
Then
\[
 \norm{E_\chi}_\infty\le C\norm{\mu_++\mu_-}_{\Carl}.       \tag{23}
\]
Indeed, write \(r=\max(|t-x|,y)\).
On \(r\le1\), the absolute scalar kernel in this integral is \(O(1)\),
and the positive dominating measure on \(r\le2^{-j}\) is \(O(A2^{-j})\).
On \(r\ge1\), since \(y\le1\), the kernel is \(O(r^{-2})\) and the
measure on \(r\le2^{j+1}\) is \(O(A2^j)\).
Both geometric series converge. The estimates are operator-order
estimates with the positive measure \(\mu_++\mu_-\).

Since \(\chi f=f\), (22)--(23) imply
\[
 f=\underbrace{\chi g_1-E_\chi}_{g}
       +S_P\bigl(\chi\mu_+-\chi\mu_-\bigr).
\]
The new measures have compact horizontal support and heights at most
\(1\), hence finite total mass by the Carleson bound.
For finite \(\N\), their positive balayages are globally in \(L_1\);
the final identity therefore holds in \(L_1\).
This proves Theorem~\ref{thm:main}.
For complex \(f\), apply the selfadjoint result to its real and imaginary
parts. Undoing the spatial normalization completes the general compact
support case.

\section*{6. General BMO and a renormalized Poisson balayage}
The compact support restriction can be removed if one normalizes the
far field. Put \(I_0=[-1/2,1/2]\), \(r(t,y)=\max(|t|,y)\), and define
\[
 \widetilde S_P\mu(x)
 =\int_{\R^2_+}\left[
 P_y(x-t)-\mathbf1_{\{r(t,y)\ge1\}}P_y(-t)
 \right]\,d\mu(t,y).                         \tag{24}
\]
Near finite points this is interpreted in local \(L_1\); at infinity
the difference is integrated in operator norm on each compact set of
\(x\)'s. Thus (24) does not subtract a potentially divergent global
constant as a separate integral.

\begin{corollary}\label{cor:general}
Under Assumption~\ref{ass:D}, for every selfadjoint
\(f\in\BMO_{cr}(\R,\N)\), there are selfadjoint \(g\in L_\infty\)
and positive operator-valued measures \(\mu_+,\mu_-\), finite on bounded
rectangles, such that
\[
 f-f_{I_0}=g+\widetilde S_P(\mu_+-\mu_-),
 \qquad
 \norm{g}_\infty+\norm{\mu_++\mu_-}_{\Carl}
 \le C(D_4)\norm{f}_{\BMO_{cr}}.                 \tag{25}
\]
The equality holds locally in \(L_1\), and modulo constants it is
independent of the reference normalization in (24).
\end{corollary}

\begin{proof}
Write \(K=\norm{f}_{\BMO_{cr}}\).
Choose a real smooth cutoff \(\eta_L\), equal to \(1\) on
\([-L,L]\), supported in \([-2L,2L]\), taking values in \([0,1]\),
and satisfying \(\norm{\eta_L'}_\infty\le C/L\).
For \(J_L=[-4L,4L]\), put
\[
 f_L=\eta_L(f-f_{J_L}),\qquad L\ge1.
\]
Then \(\norm{f_L}_{\BMO_{cr}}\le CK\), uniformly in \(L\).
Here are the details needed for the operator-valued cutoff estimate.
For an interval \(I\) of length at least \(L\),
\[
 \fint_I f_L^*f_L
 \le |I|^{-1}\int_{J_L}|f-f_{J_L}|^2
 \le 8K^2 1_\N.
\]
If \(\ell=|I|<L\) and \(I\) meets the support of \(\eta_L\),
then \(I\subset J_L\). Nested interval averaging gives
\[
 \norm{f_I-f_{J_L}}_\N
 \le CK\bigl(1+\log(L/\ell)\bigr).
\]
With \(c=\eta_L(C_I)(f_I-f_{J_L})\), split
\[
 f_L-c=\eta_L(f-f_I)
       +(\eta_L-\eta_L(C_I))(f_I-f_{J_L}).
\]
The inequality \(|a+b|^2\le2|a|^2+2|b|^2\) bounds its average square by
\[
 CK^2\left[1+(\ell/L)^2
          \bigl(1+\log(L/\ell)\bigr)^2\right]1_\N
 \le CK^2 1_\N.
\]
Variance about the actual interval mean is smaller. Intervals disjoint
from the cutoff support are trivial. The row argument is identical.

Apply Theorem~\ref{thm:main} to \(f_L\). Obtain
\[
 f_L=g_L+S_P(\mu_{L,+}-\mu_{L,-}),\qquad
 \norm{g_L}_\infty+
 \norm{\mu_{L,+}+\mu_{L,-}}_{\Carl}\le C(D_4)K.       \tag{26}
\]
Extract locally weak-star convergent subnets of the positive measures
on the closed upper half-plane, and a weak-star convergent subnet of
\(g_L\). For every compactly supported smooth \(\psi\) with
\(\int\psi=0\), and \(r=\max(|t|,y)\) sufficiently large,
\[
 |P_y\psi(t)|
 =\left|\int\bigl[P_y(t-x)-P_y(t)\bigr]\psi(x)\,dx\right|
 \le C_\psi r^{-2}.                              \tag{27}
\]
The Carleson bound gives mass \(O(Kr)\) on each dyadic annulus of radius
\(r\). Therefore the tail in pairing (26) with \(\psi\) is \(O(K/R)\)
beyond radius \(R\), uniformly in \(L\). On compact sets, \(P_y\psi(t)\)
extends continuously to \(y=0\), with value \(\psi(t)\).
Since \(f_L=f-f_{J_L}\) on the support of \(\psi\) for large \(L\),
the left side pairs exactly as \(f\).

The boundary parts of the limiting measures have bounded densities
of norm \(O(K)\), by the same argument as after (21); absorb them into
the limiting bounded function. Denote the remaining upper half-plane
measures by \(\mu_+,\mu_-\), and the bounded function by \(g_0\).
They have the required uniform bounds. We have proved
\[
 \langle f-g_0,\psi\rangle
 =\int P_y\psi(t)\,d(\mu_+-\mu_-)(t,y)
 \quad\text{whenever }\int\psi=0.                \tag{28}
\]

Formula (24) defines a local \(L_1\) function. Indeed, on a bounded
region in \((t,y)\), positivity and finite trace give local \(L_1\)
integrability of the ordinary Poisson integral. The subtracted kernel
there is bounded on \(r\ge1\). At large \(r\), the mean value theorem
bounds the difference kernel by \(C_B r^{-2}\), uniformly for \(x\)
in a compact set \(B\); the annular Carleson estimate sums this in
operator norm.
Consequently (28) implies
\[
 f=g_0+\widetilde S_P(\mu_+-\mu_-)+c
\]
for a constant operator \(c\). This constant is bounded after normalizing
\(f_{I_0}=0\). To check its bound, the average over \(I_0\) of the
renormalized balayage has norm \(O(K)\): on \(r\le2\), use
\(\fint_{I_0}P_y(x-t)\,dx\le |I_0|^{-1}\), the boundedness of
the subtracted kernel, and the \(O(K)\) mass of this region; on
\(r\ge2\), use the \(O(r^{-2})\) difference bound and annular summation.
Thus \(\norm{c}_\N\le CK\). Absorb \(c\) into \(g_0\), proving (25).
The same tail estimates justify the stated interpretation of (24).
\end{proof}

\section*{Scope and references}
The proof uses signed measures even for selfadjoint or positive \(f\).
It does not assert the stronger positivity of both \(g\) and \(\mu\).
It also does not use a dimension-independent implication from an arbitrary
positive operator-valued Carleson measure to column/row BMO; that would
be a different assertion.

\begin{itemize}
\item T. Mei, \emph{Operator valued Hardy spaces}, Mem. Amer. Math. Soc.
188 (2007), no. 881.
\href{https://arxiv.org/abs/math/0311160}{arXiv:math/0311160}.
\item M. Junge and Q. Xu, \emph{Noncommutative maximal ergodic theorems},
J. Amer. Math. Soc. 20 (2007), 385--439, especially Remark 2.2(ii).
\href{https://arxiv.org/abs/math/0505308}{arXiv:math/0505308}.
\item J. B. Garnett and P. W. Jones, \emph{BMO from dyadic BMO},
Pacific J. Math. 99 (1982), 351--371.
The random-walk argument here replaces the kernel-replacement step
of the scalar proof.
\end{itemize}
\end{document}
